\section{Численные методы расширения модели}

IF97 задаёт основной набор формул и часть обратных зависимостей,
но для полной рабочей области требуются дополнительные численные процедуры.
В проекте они используются как расширение,
сохраняющее физическую основу стандарта.

\subsection[Метод секущих для региона~3 и однофазных областей]{Метод секущих для восстановления плотности в регионе~3 и давления в однофазных областях}

В регионе~3 естественной входной парой является $(\rho, T)$.
Следовательно, при расчёте по входной паре $(p, T)$
сначала требуется восстановить плотность,
согласованную с заданными давлением и температурой.
Эта задача сводится к поиску корня уравнения
\begin{equation}
F_{\rho}(\rho) = p_{\mathrm{if97}}(\rho, T_{\mathrm{target}}) - p_{\mathrm{target}} = 0.
\end{equation}

Общий ход одномерного итерационного контура показан на рисунке~\ref{fig:secant_algorithm}.
Схема отражает последовательность выбора начального приближения,
вычисления невязки, проверки критерия остановки
и возврата к следующей итерации с учётом физических ограничений области.

\begin{figure}[htbp]
    \centering
    \begin{tikzpicture}[
        scale=0.75,
        transform shape,
        node distance=0.8cm and 1.5cm,
        auto,
        block/.style={rectangle, draw, fill=blue!10, text width=5.7cm, align=center, rounded corners, minimum height=0.8cm},
        decision/.style={diamond, draw, fill=orange!12, text width=3.2cm, align=center, inner sep=1pt, aspect=2.1},
        line/.style={draw, -Latex}
    ]
        \node[block] (start) {Входные данные: $p_{\mathrm{target}}$, $T_{\mathrm{target}}$\\или $\rho_{\mathrm{target}}$, $T_{\mathrm{target}}$};
        \node[block, below=of start] (init) {Выбор $p_{0}$ или $\rho_{0}$\\ и формирование второго приближения\\ $p_{1}=p_{0}\pm 1\%$};
        \node[block, below=of init] (calc) {Вызов вычислительного шага:\\
        $f_{0}=F(p_{0})$ или $F(\rho_{0})$\\
        $f_{1}=F(p_{1})$ или $F(\rho_{1})$};
        \node[decision, below=of calc] (tol) {$|f_{1}| \le \varepsilon$\\или\\$|f_{1}-f_{0}| \approx 0$?};
        
        \node[decision, below=of tol] (limit) {$iter \ge iter_{\max}$?};
        
        \node[block, below=of limit] (step) {Расчёт шага метода секущих:\\
        $\Delta x = -f_{1}\dfrac{x_{1}-x_{0}}{f_{1}-f_{0}}$};
        \node[block, below=of step] (clamp) {Ограничение шага и проверка\\ фазовых границ, границ региона\\ и допустимого диапазона};
        \node[block, below=of clamp] (update) {Обновление приближения:\\
        $x_{0}\leftarrow x_{1}$\\
        $x_{1}\leftarrow x_{1}+\Delta x$\\
        $iter \leftarrow iter + 1$};
        
        \node[block, fill=green!18, right=of tol] (ok) {Успешный выход:\\
        возврат согласованного состояния};
        \node[block, fill=red!15, right=of limit] (fail) {Ошибка:\\
        превышен лимит итераций\\или нарушена сходимость};

        \path[line] (start) -- (init);
        \path[line] (init) -- (calc);
        \path[line] (calc) -- (tol);
        \path[line] (tol) -- node[near start, above] {Да} (ok);
        \path[line] (tol) -- node[near start, right] {Нет} (limit);
        \path[line] (limit) -- node[near start, above] {Да} (fail);
        \path[line] (limit) -- node[near start, right] {Нет} (step);
        \path[line] (step) -- (clamp);
        \path[line] (clamp) -- (update);
        \draw[line] (update.west) -- ++(-1.4,0) |- (calc.west);
    \end{tikzpicture}
    \caption{Блок-схема одномерного итерационного контура метода секущих с контролем физических ограничений}
    \label{fig:secant_algorithm}
\end{figure}

Для одномерного восстановления используется метод секущих:
\begin{equation}
\rho_{k+1} = \rho_{k}
- F_{\rho}(\rho_{k})\,
\frac{\rho_{k}-\rho_{k-1}}
{F_{\rho}(\rho_{k})-F_{\rho}(\rho_{k-1})}.
\end{equation}
Начальные приближения $\rho_{0}$ и $\rho_{1}$
выбираются в физически допустимом интервале региона~3,
а итерации завершаются,
когда выполняется условие
\begin{equation}
|F_{\rho}(\rho_{k})| < \varepsilon_{p}
\quad \mbox{или} \quad
|\rho_{k}-\rho_{k-1}| < \varepsilon_{\rho}.
\end{equation}
После восстановления плотности
состояние рассчитывается по прямым формулам региона~3.
Опубликованные IAPWS зависимости $v(p, T)$ для региона~3
могут использоваться как опорная аппроксимация
или как источник для проверки результата~\cite{iapws_vpt_region3}.

Для однофазных регионов~1, 2 и~5
естественной входной парой является $(p, T)$.
Поэтому маршрут $(\rho, T)$
вне двухфазной области удобно трактовать симметрично:
при заданных $\rho_{\mathrm{target}}$ и $T_{\mathrm{target}}$
сначала восстанавливается давление,
после чего расчёт продолжается стандартным маршрутом $(p, T)$.
Целевая функция имеет вид
\begin{equation}
F_{p}(p) = \rho_{\mathrm{if97}}(p, T_{\mathrm{target}}) - \rho_{\mathrm{target}}.
\end{equation}

Итерационный шаг метода секущих записывается как
\begin{equation}
p_{k+1} = p_{k}
- F_{p}(p_{k})\,
\frac{p_{k}-p_{k-1}}
{F_{p}(p_{k})-F_{p}(p_{k-1})}.
\end{equation}
Критерии остановки выбираются по той же схеме,
что и для восстановления плотности:
\begin{equation}
|F_{p}(p_{k})| < \varepsilon_{\rho}
\quad \mbox{или} \quad
|p_{k}-p_{k-1}| < \varepsilon_{p}.
\end{equation}
После нахождения давления
срабатывает обычная маршрутизация IF97,
и вычисление продолжается как расчёт по входной паре $(p, T)$.
Тем самым маршрут $(\rho, T)$
остаётся встроенным в общую архитектуру модели
и не требует отдельного набора формул
для каждого однофазного региона.
